% Einheit 2 – Prozentsatz berechnen (bank/prozentrechnung/e2.jsonl, zusammenbau v0.1)
%% TODO Zweigzeile Teil 1: was hier gelernt wird (Satz in Schülersprache aus dem Einheitstitel) – Einheit 2
\einheitenkopf[e2]{Einheit 2 · Prozentsatz berechnen}
\zweigzeile{ab Kl. 7 · P10 oft}
\setcounter{aufgabe}{4}

%% TODO Titel: Ich-kann-Satz fehlt in der Bank (Platzhalter: Kettenname) – Nr. 5
%% TODO Anweisung: ein Satz über der ersten Teilaufgabe fehlt in der Bank – Nr. 5
\begin{aufgabe}{Prozentsatz}
%% TODO \rechenplatz{n}: Schrittzahl des Grundfalls fehlt in der Bank (nur bei fester Schreibform, 2.3 b)
\begin{teile}
\teil Streifen: $50$ Kinder, grau: $10$ Kinder – wie viel Prozent?

\streifenfeld{20}{$0$}{$50$ Kinder}
\teil Streifen: $40$ kg, grau: $30$ kg – wie viel Prozent?

\streifenfeld[4]{75}{$0$}{$40$ kg}
\teil Streifen: $200$ m, grau: $140$ m – wie viel Prozent?

\streifenfeld{70}{$0$}{$200$ m}
\teil Streifen: $80\,€$, grau: $40\,€$ – wie viel Prozent?

\streifenfeld[2]{50}{$0$}{$80\,€$}
\end{teile}
\end{aufgabe}

%% TODO Titel: Ich-kann-Satz fehlt in der Bank (Platzhalter: Kettenname) – Nr. 6
%% TODO Anweisung: ein Satz über der ersten Teilaufgabe fehlt in der Bank – Nr. 6
\begin{aufgabe}{Prozentsatz – weiter}
\begin{teile}
\teil $17$ von $100$ Schülern spielen ein Instrument – wie viel Prozent? \leerfeld[\%]
\teil $68$ von $100$ Losen sind Nieten – wie viel Prozent? \leerfeld[\%]
\teil $3$ von $100$ Äpfeln sind faul – wie viel Prozent? \leerfeld[\%]
\teil $45$ von $100$ Befragten fahren Rad – wie viel Prozent? \leerfeld[\%]
\teil $100$ Gäste, $81$ davon bleiben zum Essen – wie viel Prozent? \leerfeld[\%]
\teil $19$ von $50$ Kindern – wie viel Prozent? \leerfeld[\%]
\teil $7$ von $25$ Fragen richtig – wie viel Prozent? \leerfeld[\%]
\teil $13$ von $20$ Plätzen belegt – wie viel Prozent? \leerfeld[\%]
\teil $27$ von $60$ Minuten Training sind Laufen – wie viel Prozent? \leerfeld[\%]
\teil $153$ von $360$ Schülern kommen mit dem Bus – wie viel Prozent? \leerfeld[\%]
\teil $91$ von $140$ Zuschauern sind Kinder – wie viel Prozent? \leerfeld[\%]
\end{teile}
\end{aufgabe}

%% TODO Titel: Ich-kann-Satz fehlt in der Bank (Platzhalter: Kettenname) – Nr. 7
%% TODO Anweisung: ein Satz über der ersten Teilaufgabe fehlt in der Bank – Nr. 7
\begin{aufgabe}{Prozentsatz – weiter}
\begin{teile}
\teil $17$ von $24$ Kindern waren im Schwimmbad – wie viel Prozent? Runde auf eine Stelle. \leerfeld[\%]
\teil $5$ von $7$ Spielen gewonnen – wie viel Prozent? Runde auf eine Stelle. \leerfeld[\%]
\teil $38$ von $55$ Bäumen sind Kiefern – wie viel Prozent? Runde auf eine Stelle. \leerfeld[\%]
\teil Im Chor singen $14$ Mädchen und $6$ Jungen – wie viel Prozent sind Jungen? \leerfeld[\%]
\teil $18$ Siege, $4$ Unentschieden, $3$ Niederlagen – wie viel Prozent Siege? \leerfeld[\%]
\teil Eine Bäckerei verkauft $72$ Brote mit Körnern und $48$ ohne – wie viel Prozent mit Körnern? \leerfeld[\%]
\teil Jacke: vorher $80\,€$, jetzt $60\,€$ – wie viel Prozent Rabatt? \leerfeld[\%]
\teil Fahrradhelm: vorher $45\,€$, jetzt $36\,€$ – wie viel Prozent Rabatt? \leerfeld[\%]
\teil Spielkonsole: vorher $360\,€$, jetzt $306\,€$ – wie viel Prozent Rabatt? \leerfeld[\%]
\end{teile}
\end{aufgabe}

%% TODO Titel: Ich-kann-Satz fehlt in der Bank (Platzhalter: Kettenname) – Nr. 8
%% TODO Anweisung: ein Satz über der ersten Teilaufgabe fehlt in der Bank – Nr. 8
\begin{aufgabe}{Prozentsatz – weiter}
\begin{teile}
\teil Bei einer Tombola gibt es $21$ Nieten und $3$ Hauptgewinne. Wie viel Prozent der Lose sind Hauptgewinne? \hfill (P10 2018 OS) \\ \leerfeld[\%]
\teil Auf dem Blech liegen $11$ Muffins mit Blaubeeren und $4$ mit Kirschen. Wie viel Prozent sind Kirschmuffins? Runde auf eine Stelle. \hfill (P10 2018 OS) \\ \leerfeld[\%]
\teil Einwohner einer Stadt: gesamt $480\,000$; Kinder $72\,000$; Jugendliche $57\,000$; Erwachsene $264\,000$; Senioren $87\,000$. Wie viel Prozent sind Jugendliche? Runde auf eine Stelle. \hfill (P10 2023 OS)
\teil Wasser eines Haushalts pro Tag: gesamt $380$ l; Toilette $102$ l; Duschen $136$ l; Wäsche $46$ l; Rest $96$ l. Wie viel Prozent entfallen aufs Duschen? Runde auf eine Stelle. \hfill (P10 2023 OS)
\teil Von $64$ Bewerbungen wurden $47$ angenommen. Wie viel Prozent wurden abgelehnt? Runde auf eine Stelle. \hfill (P10 2015 OS)
\teil Ein Zug fuhr an $92$ Tagen, $79$-mal war er pünktlich. An wie viel Prozent der Tage war er verspätet? Runde auf eine Stelle. \hfill (P10 2015 OS)
\teil Eine Stadt hat $40\,000$ Einwohner, $1\,080$ sind in diesem Jahr zugezogen. Wie viel Prozent sind das? \hfill (P10 2014 OS) \\ \leerfeld[\%]
\teil Von $12\,500$ verkauften Konzertkarten wurden $175$ zurückgegeben. Wie viel Prozent sind das? \hfill (P10 2014 OS) \\ \leerfeld[\%]
\teil Körpergrößen der $13$ Spieler einer Mannschaft in cm: $172$, $185$, $169$, $190$, $178$, $181$, $166$, $188$, $175$, $183$, $170$, $192$, $177$. Wie viel Prozent sind größer als $180$ cm? Runde auf eine Stelle. \hfill (P10 2025 OS)
\teil Wartezeiten von $9$ Patienten in Minuten: $12$, $35$, $8$, $22$, $41$, $17$, $30$, $5$, $26$. Wie viel Prozent warteten länger als $20$ Minuten? Runde auf eine Stelle. \hfill (P10 2025 OS)
\end{teile}
\end{aufgabe}

%% TODO Titel: Ich-kann-Satz fehlt in der Bank (Platzhalter: Kettenname) – Nr. 9
%% TODO Anweisung: ein Satz über der ersten Teilaufgabe fehlt in der Bank – Nr. 9
\begin{aufgabe}{Prozentsatz – Fehler finden, Begründen}
\begin{teile}
\teil Kim, „$12$ von $48$ Kindern“: \rechnung{48 : 12 &= 4 = 400\,\%} Fehler? Wie heißt es richtig?
\teil Paul, „$18$ von $45$ Äpfeln“: \rechnung{45 : 18 &= 2{,}5 = 2{,}5\,\%} Fehler? Wie heißt es richtig?
\teil Im Korb liegen $6$ Birnen und $18$ Äpfel. Lea, Anteil der Birnen: \rechnung{6 : 18 &\approx 0{,}333 = 33{,}3\,\%} Fehler? Wie heißt es richtig?
\teil Warum teilt man beim Prozentsatz durch das Ganze und nicht durch den Teil?
\teil Ali: „$30$ von $60$ und $15$ von $30$ sind derselbe Prozentsatz.“ – Stimmt das? Begründe.
\teil Kann ein Teil mehr als $100\,\%$ seines Ganzen sein? Begründe.
\end{teile}
\end{aufgabe}

%% TODO Titel: Ich-kann-Satz fehlt in der Bank (Platzhalter: Kettenname) – Nr. 10
%% TODO Anweisung: ein Satz über der ersten Teilaufgabe fehlt in der Bank – Nr. 10
\begin{aufgabe}{Prozentsatz – Darstellungswechsel, Anwendung}
\begin{teile}
\teil Ganzes $40$ kg, Teil $10$ kg – im Streifen einzeichnen, Prozentsatz ablesen.

\streifenleer[4]
\leerfeld[\%]
\teil Ganzes $30$ m, Teil $21$ m – im Streifen einzeichnen, Prozentsatz ablesen.

\streifenleer
\leerfeld[\%]
\teil $120$ Teilnehmer, $54$ davon Frauen – im Dreisatz: wie viel Prozent?

\begin{dreisatz}{Teilnehmer}{Anteil}\dsz{120}{100\,\%}\dsp{\phantom{:0}}{\phantom{:0}}\dsz{1}{\dsleer}\dsp{\phantom{:0}}{\phantom{:0}}\dsz{54}{\dsleer}\end{dreisatz}
\teil Handyvertrag mit $20$ GB im Monat. Am 20. Tag sind $17$ GB verbraucht. Wie viel Prozent sind noch übrig? Reicht der Rest bei gleichem Verbrauch für die letzten $10$ Tage?
\teil Klassensprecherwahl in der 7c: Emma $12$ Stimmen, Noah $9$, ungültig $6$. Emma sagt: „Ich habe die Mehrheit aller Stimmen.“ Stimmt das? Runde auf eine Stelle.
\teil Eine Bäckerei wirft abends Brötchen weg: am Montag $36$ von $240$, am Dienstag $27$ von $150$. An welchem Tag war der Anteil größer?
\end{teile}
\end{aufgabe}
